\chapter{Conclusiones}
\label{sec:paper-conc}

\section{Síntesis}

GraphHunter parte de una observación simple: la caza de amenazas
es naturalmente relacional y los SIEMs la forzan a ser
tabular. La arquitectura resultante no es un producto aislado;
es la materialización de cinco disciplinas aplicadas
simultáneamente:

\begin{enumerate}
    \item \emph{Shearing layers}. Las capas de cambio rápido no
    inyectan presión sobre el núcleo. El mismo motor sirve
    hoy y servirá mañana con distintas superficies.
    \item \emph{Correctness por construcción}. El matcher SIMD tiene
    una referencia DFS y una equivalencia property-tested. La
    máquina de estados del FFI está en el sistema de tipos. Los
    scores viven en $[0,1]$ por invariante.
    \item \emph{Descomposición explicable}. El score de un path es
    la suma ponderada de cinco señales que un analista puede
    inspeccionar una por una. No hay caja negra.
    \item \emph{Degradación segura}. El GNN es opcional. La
    pérdida de conectividad a un LLM remoto no afecta la
    ingesta productiva. El sistema funciona en el primer
    día sin ajustes.
    \item \emph{Determinismo con adaptación}. El ingestor agéntico
    coloca al LLM offline, produciendo datos
    (\code{FieldConfig}), no código. El código lo produce un
    compilador determinístico. El humano aprueba, el drift
    dispara re-review.
\end{enumerate}

\section{Qué aprendimos del rediseño}

Tres lecciones merecen ser preservadas:

\begin{itemize}
    \item \textbf{El costo de una arista incorrecta es largo}. Los
    dos \code{core}$\to$\code{platform} residuales cuestan cada
    feature nueva. Invertir el tiempo en eliminarlos antes habría
    ahorrado horas acumuladas.
    \item \textbf{Los tests de equivalencia valen más que los de
    ejemplo}. La equivalencia DFS$\leftrightarrow$SIMD atrapó
    varios bugs sutiles de SIMD que ningún test unitario
    tradicional habría expuesto.
    \item \textbf{El LLM como compilador, no como traductor}. Poner
    al LLM offline y detrás de un gate humano fue la decisión que
    destrabó la ingesta adaptativa sin sacrificar determinismo.
\end{itemize}

\section{Trabajo futuro}

\subsection*{Cortes concretos, alcance de un sprint}

\begin{itemize}
    \item Eliminar las dos aristas residuales
    \code{core}$\to$\code{platform} (F2.11, F2.12).
    \item Persistir \code{mapping_library} a disco con
    compactación y retención configurable.
    \item Training pipeline reproducible para GNN por vertical.
\end{itemize}

\subsection*{Investigación, alcance mayor}

\begin{itemize}
    \item Benchmarks comparativos formales contra SIEMs comerciales
    con harness compartido.
    \item Operadores probabilísticos en el DSL (cuantificadores
    ``esta arista es inusual'') sin perder la composicionalidad.
    \item Aprendizaje en línea de pesos de scoring a partir de
    feedback del analista (tag IOC $\to$ reweighting).
    \item Multi-tenant con aislamiento por sesión y políticas de
    acceso.
\end{itemize}

\section{Cierre}

La apuesta detrás del diseño es que razonar sobre grafos de
actividad es la forma natural de investigar incidentes, y que un
sistema que respeta esa forma gana en velocidad analítica lo que
pierde en familiaridad con los paradigmas tabulares. Los
próximos seis meses de uso en producción dirán si la apuesta
paga.
